\chapter{Evaluation plan}

Chương này trình bày các chỉ số đánh giá hiệu năng (Evaluation Metrics) dành cho hệ thống Multi-Agent hiện tại và lộ trình phát triển tương lai (Future Work), tập trung vào việc tích hợp kiến trúc RAG để nâng cao khả năng suy luận và mở rộng phạm vi tự động hóa.

\section{Evaluation Metrics (Chỉ số Đánh giá)}

Khác với các hệ thống Chatbot thông thường (đánh giá dựa trên độ trôi chảy văn bản), hệ thống Security Agent được đánh giá dựa trên độ chính xác của hành động (Action Accuracy) và tính hiệu quả trong việc giảm thiểu rủi ro.

\subsection{Agentic Performance Metrics}

\textbf{Tool Selection Accuracy (TSA):}
Chỉ số này đánh giá khả năng của \textit{Scanner Agent} và \textit{Remediate Agent} trong việc lựa chọn đúng công cụ dựa trên mô tả nhiệm vụ.
\[
\text{TSA} = \frac{\text{Số lần chọn đúng Tool}}{\text{Tổng số yêu cầu gọi Tool}} \times 100
\]
Trong đó: "Chọn đúng" được định nghĩa là Tool được chọn khớp với Intent của người dùng (ví dụ: lỗi S3 thì phải chọn tool S3, không được chọn tool IAM).

\textbf{Remediation Success Rate (RSR):}
Đánh giá tỷ lệ khắc phục lỗi thành công thực tế trên môi trường AWS. Chỉ số này được đo lường tự động thông qua \textit{Rescan Agent}.
\[
\text{RSR} = \frac{| \text{Fixed Findings} |}{| \text{Attempted Fixes} |} \times 100
\]
Trong đó:
\begin{itemize}
    \item $| \text{Attempted Fixes} |$: Số lượng lỗi mà hệ thống đã cố gắng sửa.
    \item $| \text{Fixed Findings} |$: Số lượng lỗi chuyển từ trạng thái \texttt{FAIL} sang \texttt{PASS} sau khi quét lại.
\end{itemize}

\textbf{Risk Classification F1-Score:}
Đánh giá độ chính xác của \textit{Risk Evaluation Agent} khi phân loại mức độ nghiêm trọng (Critical/High/Medium/Low) so với đánh giá của chuyên gia con người (Ground Truth).
\[
F1 = 2 \times \frac{\text{Precision} \times \text{Recall}}{\text{Precision} + \text{Recall}}
\]
Mục tiêu là giảm thiểu \textit{False Negatives} (bỏ sót lỗi nghiêm trọng) và \textit{False Positives} (báo động giả gây phiền nhiễu).

\subsection{System Efficiency Metrics}

\textbf{End-to-End Latency:}
Thời gian trễ từ lúc người dùng nhập lệnh đến khi nhận được báo cáo cuối cùng.
\[
T_{total} = T_{plan} + T_{scan} + T_{analysis} + T_{remediate} + T_{report}
\]
Trong đó $T_{scan}$ thường chiếm tỷ trọng lớn nhất do phụ thuộc vào độ trễ mạng AWS API.

\textbf{Self-Correction Efficiency:}
Số vòng lặp tối thiểu cần thiết để sửa một lỗi. Một Agent tối ưu sẽ sửa lỗi ngay trong lần thử đầu tiên (1-shot) thay vì phải thử đi thử lại nhiều lần.

\section{Evaluation Methodology (Phương pháp Đánh giá)}

Chúng tôi thiết lập một môi trường kiểm thử (Sandbox Environment) với các kịch bản lỗi cố định để đánh giá hệ thống.

\subsection{Test Scenarios (Kịch bản kiểm thử)}
\begin{itemize}
    \item \textbf{Scenario A (S3 Exposure):} Tạo một S3 Bucket có quyền Public Read/Write. Kỳ vọng: Hệ thống phát hiện, đánh giá là \textit{Critical}, và có thể Block Public Access.
    \item \textbf{Scenario B (IAM Compliance):} Tạo một IAM User không có MFA. Kỳ vọng: Hệ thống phát hiện, đánh giá là \textit{High}, và gửi cảnh báo (không tự động xóa user để đảm bảo an toàn).
    \item \textbf{Scenario C (False Positive):} Tạo một tài nguyên vi phạm nhưng có gắn thẻ \texttt{Exception}. Kỳ vọng: Hệ thống nhận diện ngoại lệ và bỏ qua (Ignore).
\end{itemize}

\subsection{Human-in-the-loop Validation}
Kết quả đầu ra của \textit{Risk Evaluation Agent} sẽ được đối chiếu ngẫu nhiên (Random Sampling) bởi các chuyên gia bảo mật để tinh chỉnh Prompt Engineering.

\section{Future Development (Hướng phát triển)}

Dựa trên kết quả triển khai ban đầu (Initial Implementation), chúng tôi đề xuất lộ trình phát triển giai đoạn tiếp theo tập trung vào hai mũi nhọn: Tích hợp RAG và Mở rộng Toolset.

\subsection{Integration of Retrieval-Augmented Generation (RAG)}
Hiện tại, kiến thức bảo mật đang được "đóng băng" (frozen) trong tham số của mô hình Llama 3.2. Để hệ thống có khả năng thích ứng với các tiêu chuẩn mới mà không cần huấn luyện lại, chúng tôi sẽ tích hợp hạ tầng RAG.

\textbf{Mục tiêu:}
Cung cấp ngữ cảnh pháp lý và tiêu chuẩn tuân thủ (Compliance Standards) cho \textit{Assessment Agent}.

\textbf{Kiến trúc đề xuất:}
\begin{enumerate}
    \item \textbf{Knowledge Base:} Xây dựng cơ sở dữ liệu Vector (sử dụng ChromaDB hoặc Milvus) chứa nội dung của:
    \begin{itemize}
        \item AWS Well-Architected Framework.
        \item Tiêu chuẩn CIS Benchmark v3.0.
        \item Tài liệu nội bộ của tổ chức (Internal Policies).
    \end{itemize}
    \item \textbf{Retrieval Process:} Khi phát hiện một lỗi (ví dụ: \texttt{prowler\_check\_id: s3\_bucket\_versioning}), hệ thống sẽ truy vấn Knowledge Base để lấy ra hướng dẫn khắc phục chính xác nhất.
    \item \textbf{Benefit:} Giúp Agent giải thích lỗi "Tại sao vi phạm?" một cách thuyết phục hơn, dựa trên trích dẫn tài liệu thực tế thay vì suy luận chung chung.
\end{enumerate}

\subsection{Expansion of Remediation Toolset}
Phiên bản hiện tại tập trung chủ yếu vào dịch vụ S3. Trong tương lai, chúng tôi sẽ mở rộng thư viện \texttt{agent\_tools.py} để bao phủ toàn diện hơn:

\begin{itemize}
    \item \textbf{IAM Remediation Tools:}
    \begin{itemize}
        \item \texttt{iam\_detach\_dangerous\_policy}: Tự động gỡ bỏ quyền Administrator nếu cấp cho user không phù hợp.
        \item \texttt{iam\_force\_mfa}: Yêu cầu bắt buộc MFA cho phiên đăng nhập tiếp theo.
    \end{itemize}
    \item \textbf{EC2 Network Security:}
    \begin{itemize}
        \item \texttt{ec2\_close\_port\_22}: Tự động đóng cổng SSH (22) nếu đang mở ra 0.0.0.0/0.
        \item \texttt{ec2\_enforce\_imdsv2}: Bắt buộc sử dụng Instance Metadata Service Version 2 để chống SSRF.
    \end{itemize}
    \item \textbf{Database (RDS):}
    \begin{itemize}
        \item \texttt{rds\_enable\_encryption}: Tự động bật mã hóa lưu trữ (Storage Encryption) cho các database mới tạo.
    \end{itemize}
\end{itemize}

Việc mở rộng này sẽ chuyển đổi hệ thống từ một công cụ "S3 Auditor" thành một "Full-stack Cloud Security Guardian".